% ============================================================
% sec/05impacto.tex
%
% Impacto actual en tecnolog´ıas y organizaciones. Obligatorio. Aterrizar el tema en
%el presente: qu´e empresas, productos, servicios o sectores lo enfrentan hoy; al menos un
%caso o producto real (con nombre propio y fecha) donde el tema emergente ya cause un
%efecto medible — una brecha, un cambio de producto, una nueva norma, una inversi´on,
%un titular verificable. No basta con decir “esto es importante”; hay que mostrar d´onde y
%en qu´e ya se nota.
% ============================================================
\section{Impacto actual en tecnologías y organizaciones}
\label{sec:impacto}

El reconocimiento facial no es un tema emergente en el sentido de que
esté por llegar, sino en el sentido de que su uso ya está generando
consecuencias medibles hoy, tanto en Estados Unidos como en Costa
Rica.

En Estados Unidos, Clearview AI es el proveedor más usado por
cuerpos policiales para búsquedas de reconocimiento facial contra una
base de datos construida a partir de miles de millones de fotografías
extraídas de redes sociales. Su uso ha producido al menos dos casos
recientes de arresto injustificado con nombre propio y fecha
verificable. El primero es el caso de Angela Lipps, una mujer
de Tennessee arrestada por alguaciles federales en julio de 2025 tras
ser señalada por Clearview AI como sospechosa de un fraude bancario
cometido en Fargo, Dakota del Norte, un estado que nunca había
visitado \cite{cnn_lipps_2026}. Lipps permaneció cinco meses en
prisión hasta que sus registros bancarios demostraron que se
encontraba en Tennessee al momento de los hechos, y los cargos fueron
retirados el 24 de diciembre de 2025 \cite{cnn_lipps_2026}. El segundo
es el caso de Clarence Reid, en Jefferson Parish, Luisiana,
donde un detective utilizó Clearview AI para identificar a un
sospechoso de robo en una tienda de consignación y presentó la
identificación en el acta de arresto como proveniente de ``una fuente
creíble'', sin revelar que se trataba de un resultado de
reconocimiento facial ni corroborarlo con evidencia adicional
\cite{biometricupdate_reid_2025}. El caso terminó en un acuerdo
económico de 200{,}000 dólares en junio de 2025
\cite{biometricupdate_reid_2025}. Ambos casos muestran un patrón
repetido: el resultado del algoritmo se trata como una identificación
confirmada y no como una pista que debe verificarse, lo que conecta
directamente con la falta de una política de integridad discutida en
la Sección~\ref{sec:discusion}.

En el plano normativo, la Unión Europea ya trata la biometría como
dato de categoría especial dentro del Reglamento General de Protección
de Datos, según lo documentado por Almeida et al.
\cite{almeida_ethics_2022}. Costa Rica no cuenta con una norma
específica equivalente, pero un caso reciente muestra que la autoridad
de control local ya está aplicando los principios generales de la Ley
N.\textdegree{}~8968 a la biometría en la práctica. En mayo de 2026, la
Agencia de Protección de Datos de los Habitantes (PRODHAB) dictó la
Resolución N.\textdegree{}~029-2026-RF, que resolvió una denuncia
presentada en diciembre de 2020 contra el Tribunal Supremo de
Elecciones (TSE) por la comercialización de datos biométricos mediante
su sistema de Verificación de Identidad (VID), en operación desde
2015 \cite{lopez_prodhab_2026}. Este sistema permite a bancos y
comercios privados verificar la identidad de sus clientes mediante
cotejo de huella dactilar contra la base de datos del Registro Civil,
a cambio de un pago \cite{lopez_prodhab_2026}. PRODHAB determinó que
el TSE no contaba con habilitación legal expresa para comercializar
este tipo de datos y que existió un desvío de finalidad, ya que los
ciudadanos entregaron sus datos biométricos con un propósito
constitucional específico, la identificación para el ejercicio del
sufragio, y no para sostener una infraestructura de verificación al
servicio de terceros privados \cite{lopez_prodhab_2026}. El caso
confirma que, aunque la Ley N.\textdegree{}~8968 no menciona
explícitamente los datos biométricos, la propia autoridad de control
costarricense ya los trata como datos sensibles sujetos a los mismos
principios de finalidad y consentimiento informado que cualquier otro
dato personal, cerrando en la práctica parte del vacío normativo que
la literatura académica revisada atribuye a la ausencia de una
regulación específica sobre biometría.

En conjunto, estos hechos muestran que el problema no es hipotético.
Un proveedor comercial ya opera a gran escala dentro de cuerpos
policiales estadounidenses y ya ha producido detenciones erróneas con
consecuencias concretas sobre personas identificables
\cite{cnn_lipps_2026,biometricupdate_reid_2025}, mientras que en Costa
Rica una institución estatal ya fue sancionada por comercializar datos
biométricos sin la habilitación legal correspondiente
\cite{lopez_prodhab_2026}. Ambos casos, aunque distintos en
naturaleza, señalan la misma conclusión: la regulación específica de
la biometría sigue yendo detrás de su adopción práctica, tanto en
Estados Unidos como en Costa Rica.